% Generated from Core v2 validated Draft bytes. Do not hand-edit.
\documentclass[lualatex,a4paper,10pt]{jlreq}
\usepackage{jgaisurvey}
\addbibresource{references.bib}

\surveysetup{SP001}{Japanese Generative AI Technical Survey Special}{Thematic history / as of 2026-08-22 14:21:44 UTC}{As-of boundary: 2026-08-22 14:21:44 UTC}
\surveyeditiondescriptor{Thematic Longform Special}
\surveycoverstory{中国Generative AIの台頭}{DeepSeek・Qwen・GLM・Kimiを中心に、効率化、reasoning、agentic workload、long context、Open Weightがどのように別々の系譜から現在の競争軸へ接続したかを、2026年8月22日時点の一次資料と検証済みEvidenceから再構成する。}{Four lineages / Reasoning / Agentic / Efficiency / Serving / Open Weight / License}

\begin{document}
\surveycover
\clearpage
\section*{本号について}
\addcontentsline{toc}{section}{本号について}
本号は、中国発のGenerative AIを一つの技術潮流として平板化せず、GLM、Qwen、DeepSeek、Kimiの4系統をそれぞれのversion-specificな一次資料から追うThematic Specialである。MiniMax、Yi、Baichuanは、4系統だけでは捉えきれない並行線・歴史的bridgeとして扱い、Open Weightは技術仕様ではなくdistributionと権利条件を分けて読む。

\begin{claimboundary}[Evidence / scope boundary]
\noindent 調査・記述のas-of boundaryは2026年8月22日。後発versionの性質を過去versionへ遡及して帰属しない。\par
\noindent vendor / project / authorが報告する性能・効率値は、記載されたartifactと条件に限定して帰属する。異なるbenchmarkやserving stackを横断した勝敗表には再構成しない。\par
\noindent X / community由来の観測はDiscovery signalとして保持するが、primary corroborationのないruntime性能、普及度、license解釈を技術Evidenceへ昇格させない。\par
\noindent Open SourceとOpen Weightを同義語として扱わず、weights、code、license、commercial conditions、Model-as-a-Service条件をartifact単位で区別する。\par
\end{claimboundary}
\medskip
\tableofcontents
\clearpage
% package:SP001-P01-foundations draft-result-sha256:ae870d65136466a0ec74d8daae1098be576ab4a9cdd84ce44026d0998cce614d
\section{前史と分岐点：一つの『中国モデル』ではなく、複数の系譜が立ち上がった}
\label{pkg:SP001-P01-foundations}
\sectionkicker{THEMATIC LINEAGE 1/7 — foundations}
\noindent\textbf{GLM-130B、初期Qwen、DeepSeek LLMをそれぞれの一次資料が直接支える範囲で置くと、後年の中国発Frontier AIは単一戦略の派生ではなく、異なる設計・公開・用途の系譜が並走し始めた過程として見えてくる。} \cite{sp001sp001d001,sp001sp001d002,sp001sp001d003}\par\medskip
% block:p01-glm130b attribution:ATTRIBUTED
\noindent GLM系譜の初期座標として残るGLM-130Bは、論文上で英中バイリンガルの130Bモデルとして提示され、100B級モデルを開く試みとして位置付けられた。学習安定化やINT4 deploymentも論文の主要論点だが、ここから後年のChatGLMやGLM-4/5の設計を逆算してはいけない。重要なのは、後のGLM系譜につながる大規模モデルが、当時の一次資料でどこまで公開・運用可能性を意識していたかを、その時点の記述に限定して読むことである。 \cite{sp001sp001d001}\par\medskip
% block:p01-qwen-seed attribution:ATTRIBUTED
\noindent Qwenの最初期の技術報告は、base/chatのfamily構成とRLHFによるchat alignmentに加えて、tool use、planning、code-interpreter型のagent応用までをすでに扱っていた。したがってQwenの後年のagentic codingを突然現れた能力として描くより、早い段階からfamily/platformとtool利用を組み合わせる志向があり、それがQwen2.5以降の広い製品・モデル群へ接続したと見る方が資料に忠実である。 \cite{sp001sp001d002}\par\medskip
% block:p01-deepseek-seed attribution:ATTRIBUTED
\noindent DeepSeek LLMは、MLAやDeepSeekMoEがfamilyを特徴付ける以前のベースラインとして扱う。本Evidence passでは初期familyとrelease contextは確認できる一方、全model sizeやdistribution detailを再確定するところまでは踏み込んでいない。ここで必要なのはV2以降の効率化を過去へ投影することではなく、『後に大きく設計転換するfamilyの出発点』として明示的に分離しておくことである。 \cite{sp001sp001d003}\par\medskip
\begin{claimboundary}[Claim boundary]
この節の読解上の境界: Cross-family benchmark ranking is excluded where model versions, benchmark conditions, or evaluation methods are incompatible. / Kimi's October 2023 launch/context-window chronology remains secondary-source qualified and is not used as an unqualified foundation-date fact. / Do not project later ChatGLM or Z.ai architecture choices backward onto GLM-130B without version-specific evidence. / Later Qwen license and architecture regimes must be bound to later model versions rather than inferred from this report. / This retained evidence note does not restate every model-size or distribution detail requested by Screening; later efficiency claims are sourced to DeepSeek-V2 and subsequent artifacts. / Verify the DeepSeek LLM baseline, model sizes, release context, and open-weight distribution from primary material.: The bound primary paper verifies the early family baseline and release context, but this Evidence pass does not independently lock every requested size and distribution detail; those details are not needed for Architecture selection.
\end{claimboundary}
\clearpage
% package:SP001-P02-glm draft-result-sha256:0d805e167acaa8eca0f2632b3431d7e3261fb8c1e338bcb3944d3f9c7486878e
\section{GLM：ローカルで動く対話モデルから、All Tools、ARC、1M contextへ}
\label{pkg:SP001-P02-glm}
\sectionkicker{THEMATIC LINEAGE 2/7 — glm}
\noindent\textbf{GLM系譜では、ChatGLM-6Bのconsumer deployment、GLM-4 All Tools、GLM-4.5のARC、GLM-5.2の長文脈servingが、それぞれ別versionの資料で確認できる。公開やlicenseも含め、同じfamilyの中で戦略が段階的に変わった。} \cite{sp001sp001d005,sp001sp001d006,sp001sp001d007,sp001sp001d008}\par\medskip
% block:p02-chatglm6b attribution:ATTRIBUTED
\noindent ChatGLM-6Bは6.2Bの英中対話モデルとして、INT4で約6GBのmemoryを使うconsumer-GPU向け経路をproject側が示した。ここで技術面と同じくらい重要なのが権利の分解である。historicalなmodel-use/commercial permissionとrepository codeのlicenseは同じものではなく、『公開モデル』という一語でまとめると実際の利用条件を失う。GLMのdeveloper-facingな広がりは、動かせることと使える権利を別々に読む必要があった。 \cite{sp001sp001d005}\par\medskip
% block:p02-glm4 attribution:ATTRIBUTED
\noindent GLM-4 family reportは、GLM-4、GLM-4-Air、GLM-4-9Bへ至るversion-specificな橋を与える。報告では約10T token規模のpretrainingとmulti-stage post-trainingが説明され、All Toolsはweb browsing、Python、text-to-image、user-defined functionから必要なtoolを選ぶ能力として記述された。これはGLM系譜が単なるchat modelからtool-using systemへ軸を伸ばした転換点として読める。 \cite{sp001sp001d006}\par\medskip
% block:p02-glm45 attribution:ATTRIBUTED
\noindent GLM-4.5では、355B total / 32B activeのMoEとhybrid thinking/direct responseが論文で示され、Agentic・Reasoning・CodingをまとめたARCが明示的なfamily strategyとして前面に出た。このARC framingはGLM-4.5固有の一次資料に基づくため、他の中国モデルも同じrecipeへ収斂したと一般化すべきではない。むしろ各familyがagentic workloadへ別経路で接近したことを示す一例である。 \cite{sp001sp001d007}\par\medskip
% block:p02-glm52 attribution:ATTRIBUTED
\noindent GLM-5.2のofficial model cardでは、1M context、thinking effort control、cross-layerでsparse-attention indexを再利用するIndexShare、vLLM/SGLang deploymentが説明される。1M contextでの2.9倍per-token FLOP reductionという値も示されるが、これはvendor/model-card claimとして条件付きで扱うべき数値である。ここまで来るとGLMの競争軸はmodel qualityだけでなく、長文脈をserving可能にする設計と運用経路そのものへ広がっている。 \cite{sp001sp001d008}\par\medskip
\begin{claimboundary}[Claim boundary]
この節の読解上の境界: GLM-4.5 ARC framing is family/version-specific and is not generalized into a common Chinese-model training recipe. / Vendor efficiency figures remain source- and condition-bound rather than independent cross-family performance facts. / Model-weight/commercial permissions and repository-code licensing must not be collapsed into one license label. / Benchmark results in the paper are author evaluations and must not be merged into an incompatible cross-family ranking. / The ARC framing and reported architecture are specific to GLM-4.5 and are not evidence of a shared Chinese-model training recipe. / Efficiency numbers remain model-card/vendor claims and associated repository-code licensing may differ from model-artifact licensing.
\end{claimboundary}
\clearpage
% package:SP001-P03-qwen draft-result-sha256:0760db870f777c0deaef4d3998420e0115598401597a7cd0d2e010250ab5f6eb
\section{Qwen：広いfamily/platformを、hybrid reasoningとagentic codingへ接続する}
\label{pkg:SP001-P03-qwen}
\sectionkicker{THEMATIC LINEAGE 3/7 — qwen}
\noindent\textbf{Qwenの強みは単一の旗艦モデルではなく、specialized branchを含む広いfamily/platformを継続的に更新し、Qwen3でthinking/non-thinkingを統合し、Qwen3.6でagentic codingへ比重を移してきた点にある。} \cite{sp001sp001d009,sp001sp001d010,sp001sp001d011}\par\medskip
% block:p03-qwen25 attribution:ATTRIBUTED
\noindent Qwen2.5 technical reportは、training corpusを7Tから18T tokenへ拡大し、SFTとmulti-stage RLを組み合わせたと説明する。同時にMath、Coder、QwQ、multimodal系へつながる広いfamilyを形成しており、Qwenを一つのgeneral modelとして見るより、用途別branchを抱えるplatformとして読む根拠になる。後続のreasoning戦略は、このfamily breadthの上に構築された。 \cite{sp001sp001d009}\par\medskip
% block:p03-qwen3 attribution:ATTRIBUTED
\noindent Qwen3は0.6Bから235Bまでdense/MoEを含むfamilyを展開し、thinkingとnon-thinkingを一つの系統で扱い、thinking budgetを制御する設計をtechnical reportで提示した。reportはQwen3をApache-2.0でpublicly accessibleと位置付ける。reasoningを別モデルへ切り分けるのではなく、同じfamily内で利用形態を切り替えることが、Qwen3のdistribution strategyと相性のよい特徴になった。 \cite{sp001sp001d010}\par\medskip
% block:p03-qwen36 attribution:ATTRIBUTED
\noindent current endpointとして扱うQwen3.6-27Bのofficial model cardは、agentic coding、repository-level reasoning、過去messageをまたぐthinking contextの保持、multimodal family scope、Apache-2.0 distributionを強調する。これらはvendor claimとして読む必要があるが、Qwenの軸足が『広いopen-weight family』から『長い作業履歴を扱うcoding/agent foundation』へ連続的に伸びていることは確認できる。hosted proprietary variantとopen-weight artifactは、rightsやdistributionを論じる際には別物のままである。 \cite{sp001sp001d011}\par\medskip
\begin{claimboundary}[Claim boundary]
この節の読解上の境界: Qwen benchmark and capability statements remain tied to the named report/model card and are not normalized into a cross-family ranking. / Later Qwen architecture or licensing properties are not projected backward onto the 2023 initial Qwen release. / Hosted proprietary variants and open-weight variants must remain distinct artifacts when discussing distribution or rights. / Reported benchmark results remain evaluation-specific and are not normalized against other families. / Capability statements are vendor/model-card claims; architecture inheritance from Qwen3.5 should not be converted into independent benchmark superiority.
\end{claimboundary}
\clearpage
% package:SP001-P04-deepseek draft-result-sha256:35f64d8fd65fecb80bc2c531ce9f9c1a6d2f9371a246c5fd52085272385d70df
\section{DeepSeek：効率をarchitectureで作り、reasoning RLとagentic servingへ伸ばす}
\label{pkg:SP001-P04-deepseek}
\sectionkicker{THEMATIC LINEAGE 4/7 — deepseek}
\noindent\textbf{DeepSeekの転換は、V2のMLA/DeepSeekMoE、R1/R1-Zeroのreasoning RL、V4-Flash-0731のagentic servingという別々の節目で追える。benchmark順位ではなく、各versionで何が設計・学習・deploymentの焦点になったかを見る。} \cite{sp001sp001d012,sp001sp001d013,sp001sp001d014}\par\medskip
% block:p04-v2 attribution:ATTRIBUTED
\noindent DeepSeek-V2は236B total / 21B activeのMoEで、Multi-head Latent Attention（MLA）とDeepSeekMoEをfamily-definingな効率化として導入した。128K contextとともに、paperはKV-cacheやgeneration throughputの改善を報告する。ただし数値比較はpaperのsetupに依存する。ここで重要なのは、DeepSeekがfrontier competitionへ入る際、parameter countの増大だけでなくmemory/serving costをarchitectureの中心課題に置いたことである。 \cite{sp001sp001d012}\par\medskip
% block:p04-r1 attribution:ATTRIBUTED
\noindent DeepSeek-R1 paperは、R1-Zeroをpreliminary SFTなしのlarge-scale RLとして、R1をcold-start dataとmulti-stage trainingを加えた別の経路として説明する。またQwen/Llamaをbaseにしたdistilled variantsもreleaseされた。native R1/R1-Zeroとdistillを一つのmodel subjectとして扱わないことが重要であり、reasoningの成果をどのtraining pathとmodel lineageへ帰属させるかを分離する必要がある。 \cite{sp001sp001d013}\par\medskip
% block:p04-v4 attribution:ATTRIBUTED
\noindent V4-Flash-0731のofficial model cardから直接確認できるcurrent endpointは、agentic capabilityの強化とvLLM、DSpark speculative decodingを用いるdeployment supportである。一方、この0731 artifactだけではbroader V4 previewのarchitecture全体を再確定できない。したがって本号では、V2で確立した効率architecture、R1のreasoning post-training、V4-Flashのagentic servingという確認可能な線だけをつなぎ、別sourceのV4仕様を補完的に持ち込まない。 \cite{sp001sp001d014}\par\medskip
\begin{claimboundary}[Claim boundary]
この節の読解上の境界: Paper-reported efficiency comparisons are setup-bound and are not generalized across unrelated serving stacks. / The broader DeepSeek V4 architecture is not silently imported into the V4-Flash-0731 Evidence boundary. / Paper-reported efficiency comparisons must retain the stated comparison setup and must not be generalized across unrelated serving stacks. / Native R1/R1-Zero and distilled Qwen/Llama-based variants are different model subjects and must remain separately attributed. / The bound 0731 update is sufficient for agentic/serving endpoint evidence but does not independently restate every architecture detail of the broader V4 preview. / Verify V4-Flash-0731 context architecture, agentic positioning, and official vLLM/DSpark deployment support.: The official 0731 model card verifies agentic positioning and vLLM/DSpark deployment support; full context-architecture verification is retained as a bounded limitation rather than inferred from separate V4 sources.
\end{claimboundary}
\clearpage
% package:SP001-P05-kimi draft-result-sha256:73b3a6ad971550d2698af6be43d38b98f3d1ff23481e36af266c9dcc13b32e2f
\section{Kimi：long contextを起点に、reasoning RL、open agentic MoE、native multimodalへ}
\label{pkg:SP001-P05-kimi}
\sectionkicker{THEMATIC LINEAGE 5/7 — kimi}
\noindent\textbf{Kimiの2023年launch chronologyは二次資料の制約を残す一方、2025年以降はk1.5、K2、K3の一次・official materialによって、long-context researchからagentic・multimodal frontierへ伸びる技術線を確認できる。} \cite{sp001sp001d015,sp001sp001d016,sp001sp001d017}\par\medskip
% block:p05-k15 attribution:ATTRIBUTED
\noindent Kimi k1.5 paperは、multimodal reasoning modelをscaled reinforcement learningで訓練し、long-context scaling、policy optimization、infrastructure optimization、long-to-short transferを主要要素として説明する。2025年1月のchronologyはDeepSeek-R1と近接するが、両paperが開示するtechniqueは同じではない。この近さから模倣や単純な先後関係を推論せず、Kimi側のlong-context RL研究として独立に位置付ける。 \cite{sp001sp001d015}\par\medskip
% block:p05-k2 attribution:ATTRIBUTED
\noindent Kimi K2では、1T total / 32B activeのMoE、MuonClip optimization、15.5T-token pretrainingがpaperで説明され、agentic data synthesisとreal/synthetic environmentをまたぐjoint RLへ踏み込んだ。k1.5がlong-context reasoningの研究線を明確にしたのに対し、K2はそれをopen agentic modelのdistributionへ接続する節目として扱える。 \cite{sp001sp001d016}\par\medskip
% block:p05-k3 attribution:ATTRIBUTED
\noindent Kimi K3のofficial materialは、2.8T total / 104B active、native multimodal、1M context、KDAとgated MLAを備え、long-horizon codingやagentic knowledge workを狙うcurrent endpointとしてmodelを位置付ける。これらはofficial/vendor materialに基づく記述であり、license条件はmodel cardから推測しない。K3のcommercial条件は別のexact LICENSE artifactに委ねることで、技術仕様と権利条件を分離して扱う。 \cite{sp001sp001d017}\par\medskip
\begin{claimboundary}[Claim boundary]
この節の読解上の境界: Kimi's October 2023 launch/context-window chronology remains secondary-source qualified until a stronger first-party archival source is added. / Kimi k1.5 and DeepSeek-R1 are treated through their disclosed techniques and chronology without unsupported imitation or priority claims. / The contemporaneous k1.5 and DeepSeek-R1 papers use different disclosed techniques; priority or simple imitation claims are not supported. / Performance claims remain paper evaluations and should not be normalized against incompatible benchmarks from other families. / Detailed commercial/license conditions are verified separately against the exact Kimi K3 LICENSE artifact in SP001-D021. / Verify Kimi K3 architecture, parameter structure, KDA/Gated MLA, multimodal/agentic claims, distribution, and applicable license boundary from official material.: The official model material verifies architecture, parameter structure, multimodal/agentic positioning and distribution; exact legal conditions are intentionally delegated to the separate exact-license Evidence card rather than inferred here.
\end{claimboundary}
\clearpage
% package:SP001-P06-parallel draft-result-sha256:d036ff68cf50cdc5615a282aab5182652c7530c5897b8ba0dd5fafffb34a25f3
\section{主要4系統の外側：MiniMax、Yi、Baichuanが示す競争と分化}
\label{pkg:SP001-P06-parallel}
\sectionkicker{THEMATIC LINEAGE 6/7 — parallel}
\noindent\textbf{中国発AIの競争史を主要4系統だけで閉じると、agentic frontierの別経路、2023-24年のopen-weight/local wave、domain specializationへの分化を見落とす。MiniMax、Yi、Baichuanはその補助線になる。} \cite{sp001sp001d018,sp001sp001d019,sp001sp001d020}\par\medskip
% block:p06-minimax attribution:ATTRIBUTED
\noindent MiniMax M2.7のofficial announcementは、Agent Teams、skills、dynamic tool search、さらにmodelが自身のdevelopment workflowの一部へ参加するという方向を打ち出す。これは2026年のagentic frontierを主要4familyだけで説明できないことを示すmaterial competing branchである。ただし、このannouncement単体からopen-weight serving integrationやexact licenseまで確定することはせず、確認できるagentic positioningに役割を限定する。 \cite{sp001sp001d018}\par\medskip
% block:p06-yi attribution:ATTRIBUTED
\noindent Yi-1.5 repositoryは34B/9B/6B variants、Ollamaを含むlocal use、fine-tuning ecosystem integration、Apache-2.0 code/weightsを示す。これは2023-24年に、中国発open-weight modelがdeveloperのlocal environmentへ入り込む経路が複数存在したことを示すsupporting bridgeである。一方、この2024 repositoryだけから2026年にも独立したgeneral-frontier lineが続いているとは結論しない。 \cite{sp001sp001d019}\par\medskip
% block:p06-baichuan attribution:ATTRIBUTED
\noindent Baichuan-M3-235Bのofficial model cardは、medical-enhanced clinical inquiryへのspecialization、Qwen3-MoE ancestry/acknowledgement、Apache-2.0 licenseを示す。ここではBaichuanを現在のco-equal general frontierとして数えるのではなく、初期の独立brandがdomain-specializedかつ他family由来のbaseを使う方向へ分化しうるecosystem例として置く。これは競争が『familyの数』だけでなく、派生・専門化の形でも進むことを示している。 \cite{sp001sp001d020}\par\medskip
\begin{claimboundary}[Claim boundary]
この節の読解上の境界: MiniMax is a visible competing branch, not a fifth co-equal central family in this thematic architecture. / X runtime-adoption and enterprise-origin-policy observations remain HOLD and cannot support prevalence or performance claims. / The bound news source does not itself establish every requested open-weight serving integration or exact license term; those were observed in a separate official model card during Discovery but are outside this one-source Evidence task. / Verify MiniMax M2.7 Agent Teams, skills/tool-search claims, model-assisted development positioning, open-weight distribution, serving support, and license from official sources.: The official announcement verifies Agent Teams, skills/tool search and model-assisted-development positioning; open-weight serving and license details are not promoted from a different source into this single-source card. / The Yi-1.5 repository alone cannot prove the absence of a later independent general Yi frontier endpoint; that negative conclusion remains a Discovery-level catalog observation. / Verify Yi-1.5 release, open/local deployment, fine-tuning integrations, licensing, and whether an official later general Yi endpoint exists within the as-of boundary.: The repository verifies Yi-1.5 release/deployment/fine-tuning/licensing; the absence of a later co-equal general endpoint is not established by this source alone and remains contextual. / The historical transition from the earlier independent Baichuan2 foundation-model line requires separate earlier evidence and is not established by the M3 model card alone. / Verify Baichuan M3 medical specialization, Qwen3-MoE ancestry, licensing/deployment, and the historical transition from the earlier independent Baichuan line.: The M3 model card verifies medical specialization, Qwen3-MoE ancestry/acknowledgement and Apache-2.0 licensing; the earlier-to-current historical transition is only partially represented by this source.
\end{claimboundary}
\clearpage
% package:SP001-P07-open-weight draft-result-sha256:6999f0e7d1bd3616c110a54222e0aa1ff08273c82c34f326572e6d1f59915f9f
\section{Open Weightは一枚岩ではない：配布されることと、同じ権利で使えることは別問題}
\label{pkg:SP001-P07-open-weight}
\sectionkicker{THEMATIC LINEAGE 7/7 — weight}
\noindent\textbf{公開weightsはdeveloper adoptionを加速しうるが、code、weights、commercial use、attribution、Model-as-a-Service条件はartifactごとに異なる。Kimi K3のexact licenseとYi/BaichuanのApache-2.0 artifactsを並べると、その差が具体的になる。} \cite{sp001sp001d021,sp001sp001d019,sp001sp001d020}\par\medskip
% block:p07-k3-license attribution:FACTUAL
\noindent Kimi K3 Licenseは、covered software/weightsについてuse、copy、modify、distribute、sublicense、sell、deploy、fine-tuneを含む広いgrantを置く一方、定義されたModel-as-a-Service businessが連続12か月のaggregate revenueで2000万米ドルを超える場合にはcommercial use前のseparate agreementを求める。また対象となる大規模commercial product/serviceには、100M MAUまたは月間revenue 2000万米ドル超を条件とする『Kimi K3』表示条項があり、license本文には除外も記されている。これはK3という特定artifactの条件であり、Kimi family全体へ遡及させない。 \cite{sp001sp001d021}\par\medskip
% block:p07-apache-comparison attribution:ATTRIBUTED
\noindent 比較対象としてYi-1.5はrepository上でApache-2.0 code/weightsとlocal/fine-tuning integrationsを示し、Baichuan-M3もofficial model card上でApache-2.0を掲げる。これらをKimi K3と並べる目的は『どれがよりopenか』を単純順位付けすることではない。公開weights、repository code、商用条件、attribution、派生物の扱いをartifact単位で分けて読む必要があることを示すためである。 \cite{sp001sp001d019,sp001sp001d020}\par\medskip
% block:p07-boundary attribution:MIXED
\noindent 本号でのlicense整理はtechnical distributionを理解するための境界整理であり、法的助言ではない。X上のlicense-friction discussionはHOLDのままで、infringement、derivative provenance、適用version、third-party complianceの判断根拠には用いない。『Open Source』と『Open Weight』も同義語として扱わず、確認できるweights/code/license条件だけを個別に記述する。 \cite{sp001sp001d021}\par\medskip
\begin{claimboundary}[Claim boundary]
この節の読解上の境界: The Kimi K3 License summary is artifact-specific and is not a legal opinion. / X license-friction discussion remains HOLD and cannot establish infringement, derivative provenance, applicable version, or legal interpretation. / This Evidence card summarizes the named 2026 Kimi K3 License only; it is not a legal opinion and does not establish infringement or compliance by third parties. / The Yi-1.5 repository alone cannot prove the absence of a later independent general Yi frontier endpoint; that negative conclusion remains a Discovery-level catalog observation. / Verify Yi-1.5 release, open/local deployment, fine-tuning integrations, licensing, and whether an official later general Yi endpoint exists within the as-of boundary.: The repository verifies Yi-1.5 release/deployment/fine-tuning/licensing; the absence of a later co-equal general endpoint is not established by this source alone and remains contextual. / The historical transition from the earlier independent Baichuan2 foundation-model line requires separate earlier evidence and is not established by the M3 model card alone. / Verify Baichuan M3 medical specialization, Qwen3-MoE ancestry, licensing/deployment, and the historical transition from the earlier independent Baichuan line.: The M3 model card verifies medical specialization, Qwen3-MoE ancestry/acknowledgement and Apache-2.0 licensing; the earlier-to-current historical transition is only partially represented by this source.
\end{claimboundary}
\clearpage
\section{この号の総括}
\label{sec:issue-summary}
\sectionkicker{ISSUE SYNTHESIS}
\noindent 本号を通して見えるのは、『中国モデル』という単一のarchitectureや学習recipeではなく、異なる出発点を持つ複数の系譜である。GLMは公開・ローカル展開からtool use、ARC、長文脈servingへ、Qwenは広いfamily / platformからhybrid reasoningとagentic codingへ、DeepSeekはMLA / MoEによる効率設計からreasoning RLとagentic servingへ、Kimiはlong contextを起点にreasoning RL、open agentic MoE、native multimodalへ進んだ。相互に近い課題へ向かっていても、その技術的な道筋を一つにまとめることはできない。\par\medskip
\noindent 一方で、2025年から2026年にかけて競争課題そのものは重なってきた。reasoningをどのように利用形態へ組み込むか、codingやtool useを長時間のagentic workloadへどう接続するか、long contextを実際にserveできる効率へ落とすか、そしてweightsを公開したモデルをdeveloperがどのruntimeやdistribution経路で使えるようにするかである。収斂しているのは一つの実装方法ではなく、frontier modelに求められる問題設定の側だと捉える方が、本号で確認したEvidenceには整合的である。\par\medskip
\noindent MiniMax、Yi、Baichuanを併記すると、4主要系統だけで中国のopen-weight史を閉じることもできない。MiniMaxはagentic frontierの並行線として、Yiは2023〜2024年のopen distributionを示すbridgeとして、Baichuanは初期の独立系譜から現在のspecialized / derivativeな位置へ移る例として機能する。そしてOpen Weightという言葉も一つの制度ではない。weightsの取得可能性、repository code、再配布、商用利用、表示義務、Model-as-a-Service条件はartifactごとに異なり、技術的な『開放性』と法的・商用上の条件を分離して確認する必要がある。\par\medskip
\noindent なお、本号には意図的に残した境界がある。Kimiの2023年初期chronologyは一次資料で完全には閉じておらず、二次資料としての資格を外していない。X / communityのruntime採用やlicense frictionも、技術仕様や普及度を確定するEvidenceには使っていない。また、各社・各論文のbenchmark値を異なる条件のまま横断順位へ変換していない。したがって本号が示すのは『どのモデルが絶対に最強か』ではなく、2026年8月22日までに各系譜が何を競争課題として選び、どの公開・serving・権利境界を伴ってfrontierへ到達したか、という技術史である。\par\medskip
\clearpage
\printbibliography[title={References / Source Notes}]
\end{document}
